\color{unidarkgreen}{\subsection{Защита от непереключения фаз (ЗНФ)}}
\color{black}

\begin{enumerate}[label=\arabic{section}.\arabic{subsection}.\arabic*,  align=left, labelsep=4pt, leftmargin=0pt, itemindent=57pt]
\item
Защита от непереключения фаз (ЗНФ) предназначена для обнаружения расхождения полюсов
выключателя с пофазным приводом при его включении или отключении. 

\begin{figure}[H]
\centering
\includegraphics[width=1\textwidth,height=1\textheight,keepaspectratio]{\fbpath/761.3441 ЗНФ/img/ZNF.pdf}
\captionof{figure}{Функциональная схема ЗНФ В1(2)}\label{pict_znf:ZNF}
\end{figure}

\item
Ввод в работу ЗНФ выключателя В1(2) осуществляется с помощью программных переключателей \textbf{«ЗНФ
В1(2)»}. Факт расхождения полюсов может устанавливаться одним из следующих способов:

\begin{enumerate}[label=\arabic*)]
\item на основании внешнего сигнала от сборки блок-контактов выключателя;
\item функцией фиксации положения выключателя (ФПВ), которая пофазно контролирует состояние блок-контактов выключателя и в случае расхождения полюсов формирует сигнал \textit{«Внутр. пуск. ЗНФ В1(2)»}.
\end{enumerate}

\item
В зависимости от используемого способа соответствующий сигнал должен назначаться на
конфигурируемый вход \textit{«Пуск ЗНФ В1(2)»}.

\item
Действие на отключение выключателя В1(2) с запретом АПВ осуществляется с выдержками времени
\textbf{«Тср ЗНФ В1(2)»}. Задаваемая выдержка времени предназначена для отстройки от разновременности
переключения блок-контактов выключателя.

\item
Предусмотрена возможность блокировки ЗНФ В1(2) от внешнего сигнала \textit{«Блок. ЗНФ В1(2)»} на время,
которое задается уставкой \textbf{«Тдеблок ЗНФ В1(2)»}.

\item
Оперативный ввод/вывод ЗНФ В1(2) производится с помощью БУ (блоков управления функциями) \textit{«Вывод ЗНФ В1(2)»}. 

\small
\begin{longtable}{|>{\centering\arraybackslash}m{5.3cm}|>{\centering\arraybackslash}m{3.3cm}|>{\centering\arraybackslash}m{4.2cm}|>{\centering\arraybackslash}m{1.8cm}|>{\centering\arraybackslash}m{1cm}|}
\caption{Перечень уставочных параметров ЗНФ\hfill\vspace{-0.5\baselineskip}}\label{znf:tbl1}\\ 
\hline
\rowcolor{gray!30}
Параметр & Условное обозначение на схеме & Значение/ Диапазон & Единица измерения & Шаг \\ 
\hline
\endfirsthead
\caption*{\hspace{3pt}\emph{Продолжение таблицы \ref{znf:tbl1}\hfill\vspace{-0.5\baselineskip}}} \\ % сделано по ГОСТ 2.105 п.6.8.7
\hline
\rowcolor{gray!30}
Параметр & Условное обозначение на схеме & Значение/ Диапазон & Единица измерения & Шаг \\ 
%\hline
\endhead

%===m>ЗНФ|-
%===m
\end{longtable}
\normalsize

\end{enumerate}